\cotefile{156-9}
\scanpage{10}
\mspage{8}
% source: top, two-line display
\[
\left\{\begin{aligned}
I_{<a}\overset{\unsure{\circ|0}}{\ll}I_{<b}&\Longleftrightarrow a\leq b\\
X\overset{\circ}{\ll}I_{<a}&\deleted{\unsure{\&}}\Longleftrightarrow X\deleted{\unsure{<}}\leq a
\end{aligned}\right.
\qquad \deleted{\unsure{\text{et}}}\quad\text{(cf déf\ldots)}
\]
% source: upper-left margin
% margin note written sideways; slanted
\margin{On \uuline{\unsure{a plus|de plus}} $X\overset{\circ}{\ll}Y\Longleftrightarrow |X|\subset|Y|$ et $X\unsure{\overline{\triangleleft}|\overline{\triangle}}Y$ pour $X,Y\in\unsure{\widehat M|\widehat{\mathscr M}}$ et \unsure{remarquer} \deleted{\unsure{(*)}}\unsure{?} C'est vrai \unsure{seult|seulement} si $\unsure{L|\mathscr L}$ n'a ni plus petit ni plus grand élément.}
% source: top, continuation below display
(où $X\in M$), et \unsure{aussi} pour \deleted{$\unsure{I_{<a}}$} les relations \deleted{\unsure{des semi}} impliquant $I_{>a}$ \unsure{intervalles} \unsure{infinis|définis} symétriquement.

On note.

% source: upper-middle, a)
a) On n'a jamais une relation
\[
I_{<a}\overset{\circ}{\ll}X\qquad\text{pour }X\in M
\]
même si $L$ a un plus petit él. $\unsure{a_0|a_o}$.

% source: middle, example
Par exemple, \deleted{\unsure{que nous}} on n'a \uline{pas} $I_{<a}\overset{\circ}{\ll}I_{\omega_0,a}$, \deleted{car (\unsure{la même heuristique},} \unsure{donc pas} $I_{\omega_0,a}\overset{\circ}{\ll}I_{<a}$. En fait, $I_{\omega_0,a}$ doit être vu (en tant que définissant une \unsure{multiplicité}\ldots) comme une \uline{subdivision} effective de $I_{<a}$.

(Si on avait les deux relations à la fois, on \deleted{aurait}\inserted{devrait avoir} $I_{<a}=I_{\omega_0,a}$, ce qui n'est pas conforme aux intuitions \unsure{présidant} aux \unsure{introductions} de $I_{<a}$.
% source: small added paragraph below preceding lines
\inserted{\unsure{Bien sûr}, \unsure{on aura|on n'aura} $\deleted{\unsure{\not\ll}}I_{<a}\ll I_{\omega_0,a}$, mais la relation inverse \deleted{\unsure{seule}} $I_{\omega_0,a}\ll I_{<a}$, est vraie malgré que $|I_{\omega_0,a}|=|I_{<a}|$.}

% source: lower-middle, b)
b) On n'a jamais une relation
\[
I_{<a}\overset{\circ}{\ll}I_{>b},\quad\text{ou l'inverse.}
\]
% source: lower-left margin
% slanted
\margin{\unsure{Ne plus|Ne pas}\unsure{?} sur $\widehat M$}
% source: bottom, cancelled definition
Relation d'ordre $\leq$ sur $\widehat M$. On \unsure{pose}, \inserted{\unsure{encore},} \deleted{si $\unsure{X,Y\in}\ X,Y\in\widehat M$} \unsure{\text{Ord}} $\leq:\widehat M$ \unsure{par}
\[
\deleted{X\leq Y\Longleftrightarrow\forall x\in\delta X,\ \forall y\in\delta Y,\quad x\leq y}
\]
\deleted{$\unsure{I_{<}}$}
\deleted{(on écrit aussi $\delta X\leq\delta Y$ relation qui a un sens pour deux parties $A,B$ quelc de $M$: $A\leq B$)}
% archive sheet number: 9
\scanpage{11}
\mspage{7}
% source: top, cancelled heading and replacement definition
\deleted{Relation $\unsure{\mid\circ\mid|\mid\bullet\mid}$ sur $\unsure{\widehat M|\widehat{\mathscr M}}$}
\[
X\unsure{<|\leq}Y\Longleftrightarrow\inserted{X\ne Y\text{ et }}|X|\leq|Y|
\]
On trouve
\[
\deleted{X\mid\circ\mid Y\Longleftrightarrow X<Y\text{ ou }Y<X\quad\text{i.e. }\{X,Y\}\in\unsure{\operatorname{Drap}_2}(\widehat M,\leq\unsure{\}|)}}
\]
% source: upper part, order definition
\deleted{$\unsure{M|A}$}
\[
X\leq Y\Longleftrightarrow
\left\{\begin{aligned}
&1^\circ)\ \delta X\leq\delta Y\\
&2^\circ)\ \text{Si }X\text{ ou }Y\text{ ou les deux de la forme }I_{<a},I_{>\unsure{b|a}}\\
&\qquad\text{on exige }\deleted{\text{\unsure{pour eux}}}\text{ ceci:}\\
&\qquad a)\ \text{Si }\deleted{X}Y=I_{<a},\ \deleted{\text{alors }X=Y}\text{ alors }X=Y\\
&\qquad b)\ \text{Si }X=I_{>a},\ \text{alors }Y=X
\end{aligned}\right.
\]
\deleted{\unsure{et Si} $\unsure{X\leq I}$}

% source: upper-middle, first NB
\uline{NB} On a dans tous les cas
\[
\begin{aligned}
X\leq Y&\Longrightarrow |X|\leq|Y|\\
X<Y&\Longrightarrow |X|<|Y|
\end{aligned}
\]

% source: middle, second NB
\uline{NB}\deleted{\unsure{On}} On n'a jamais $I_{>a}\leq I_{<b}$ à cause du $2^\circ$ a) (ou b)) ci-dessus. On a, bien sûr,
\[
I_{<a}\leq I_{\deleted{<}>b}\Longleftrightarrow a\leq b.
\]

% source: lower-middle, definition of barred-dot relation
On pose \deleted{$\unsure{\varphi|\phi}$}, pour $\unsure{X,Y|x,y}\in\widehat M$
\[
X\mid\circ\mid Y\Longleftrightarrow X<Y\text{ ou }Y<X\quad\text{i.e. }\{X,Y\}\in\operatorname{Drap}_2(\widehat M,\leq\}.
\]

% source: bottom, axiom verification
Vérification des axiomes. \unsure{Mog|Mo$\rho$} 1 OK, Mog 2 OK, Mog 3 OK car
\[
X\deleted{\unsure{\leq}}<Y,\ X'\overset{\unsure{\circ|0}}{\ll}X,\ Y'\overset{\circ}{\ll}Y\Longrightarrow X'<Y'.
\]
\deleted{\unsure{si on}}\inserted{\unsure{car}}
\diagram{\begin{tikzcd}[ampersand replacement=\&,column sep=small,row sep=small]
\underbrace{X<Y,\ X'\ll X,\ Y'\ll Y}\arrow[r,Rightarrow]\arrow[dd,Rightarrow,bend right=25]\& X'\leq Y'\\
\&\deleted{\unsure{\longleftarrow X'\equiv}}\\
X'<Y'\text{ ou }X'=\unsure{\operatorname{ex}X|\operatorname{ec}X|e_X}=Y'=\unsure{\operatorname{or}Y}\&
\end{tikzcd}}
\uline{Mog 4} clair.
% archive sheet number: 10
\scanpage{12}
\mspage{8}
% source: top, cancelled first lemma
\deleted{Proposition}

\deleted{\uline{Lemme} a) Soit $a\in\unsure{L|\mathscr L}$ un élément \unsure{non terminal|non minimal}, en sorte que $I_{<a}$ est défini. Alors}
\[
\deleted{\mathrm{cos}\!\!\unsure{-}\!{\unsure{w|W|\psi}}^{\unsure{\circ|0}}(I_{<a})=\{a\}\cup\operatorname{Omb}^{\circ}(I_{>a})}
\]

% source: upper-middle, new lemma
\uline{Lemme} $\unsure{M|\mathscr M}$ est très dense dans $\widehat M$, i.e. pour $X,Y\in\widehat M$, pour qu'on ait $X\unsure{\mid\circ\mid|\mid\bullet\mid}Y$, il faut que pour tout $X'\overset{\circ}{\ll}X$, $Y'\overset{\circ}{\ll}Y$, avec $X',Y'\in M$, on ait $X'\mid\circ\mid Y'$.
% source: left margin adjacent to lemma and proof
% margin note written sideways; slanted
\margin{\uline{NB} \unsure{Il suffit} de \unsure{le|la} vérifier en prenant $X'\overset{\circ}{\ll}X$, et \unsure{pareil} $Y$.}

% source: middle, first proof paragraph
\deleted{Supp.} Il n'y a de question \deleted{que \unsure{si jamais}}\inserted{si $X$ est un} \deleted{\unsure{au moins un} des deux élém $X,Y$ est un} $I_{<a}$ ou un $I_{>b}$. Supposons p.ex.
\[
X=I_{<a},
\]
\deleted{et supposons d'abord $Y\in M$,} Alors \deleted{\unsure{X est}} $X\mid\circ\mid Y$ sss $I_{<a}<Y$ i.e. $a\leq Y$, \deleted{et on a pour $X\overset{\circ}{\ll}I_{<a}$,}\inserted{montrons qu'il suffit} $\forall X'\overset{\circ}{\ll}I_{<a}$, on ait $X'\mid\circ\mid Y$.

\deleted{\unsure{c.à.d.} $X\unsure{\leq|<}a$, \unsure{alors}}

% source: lower-middle, choice of X prime
Prenons d'abord $X'=\deleted{[a',a]}I_{\inserted{a',a}}$, pour $a'<a$ (tel $a'$ existe, car $a$ non minimal). On aura donc $Y\mid\circ\mid I_{a',a}$ donc $Y\geq a$ \uline{ou} $Y\leq a'$. Mais si on avait $Y\leq a'<a$, on aurait $Y\overset{\circ}{\ll}X=I_{<a}$, et prenant $X'=Y$, on trouve $Y\mid\circ\mid Y$, absurde.

% source: bottom, cancelled continuation
\deleted{\unsure{Si} On prouve de \unsure{même}, pour $X'$ fixé, que $X'\mid\circ\mid Y$ sss $\unsure{Y}$}

% source: bottom, corollary
\uline{Cor.} L'ens\inserted{(Collection ou cl.)} des supports $\Sigma\widehat M$ est can. isom à l'ens des supports $\Sigma M$.
% archive sheet number: 11
\scanpage{13}
\mspage{9}
% source: top, continuation
Mais c'est à l'ens des supports \uline{pour $\unsure{M|\mathscr M}$} qu'on s'intéresse en fait!

% source: upper part, figures and polygons
Nous allons à présent travailler avec des figures et polygones finis de $\widehat M$ (ou $\widehat M$-polygones, ou polygones \unsure{généralisés}), i.e. les parties finies tot. \unsure{ordonnées} (ou drapeaux) de $\widehat M,\leq$. Pour une telle partie $\unsure{\Phi|\varphi}$, on introduit
\[
T^{\Phi}=\delta\Phi\in\operatorname{Drap}^{\unsure{\times|*}}(\unsure{L|\mathscr L}),
\]
qui est vide ssi $\Phi$ est vide.
% source: left margin beside definition of T
% slanted
\margin{\unsure{On va par la} suite se borner aux non vides.}

% source: middle, interstitial strata
Ainsi on introduit les \uline{strates interstitielles} (\unsure{généralisées|généralisation}) d'un drapeau
\[
T=\{t_{\unsure{1|l}}<t_2<\cdots<t_n\}
\]
de $L$, ainsi: \deleted{\unsure{Si on prend}}
\[
[I_{<t_1}],\ t_1,\ I_{t_1,t_2},\ t_2,\ I_{t_2,t_3},\ \cdots,\ t_n,[I_{>t_n}]
\]
% source: lower-middle, qualifying sentence and insertion
où le terme $I_{<t_1}$ ne figure que si $t_1$ non minimal\inserted{(sinon, il n'existe \unsure{même} pas)}, et le terme $I_{>t_n}$ ne figure que si $t_n$ est non maximal (\unsure{idem}\ldots).

% source: lower, largest polygon
Ce sont les éléments de le plus grand $\widehat M$-polygone possible $\widehat\Phi$, telle que $\delta\widehat\Phi=T$.

% source: bottom, data defining Phi
La donnée de \inserted{\unsure{le} $\widehat M$-polygone} $\Phi\ne\varnothing$ revient à celle du drapeau $T$ plus un ensemble fini\inserted{$\Phi$,} des strates interstitielles (\unsure{généralisées}) ayant la propriété suivante:
\[
\forall t\in T,\quad\text{on a}\quad t\in\Phi,\quad\uline{\text{ou}}\quad t\triangleleft X\quad\text{avec }X\in\Phi.
\]
% archive sheet number: 12
\scanpage{14}
\mspage{10}
% source: top, complementary polygon
Ceci posé, on définit le \uline{polygone complémentaire} ${\unsure{\Phi|\varphi}}^{*}$ de $\Phi$ comme
\[
\begin{aligned}
\Phi^*&=\widehat\Phi\setminus\Phi=\text{ens des }\deleted{\text{intervalles}}\inserted{\text{strates}}\\
&\qquad\text{interstitielles généralisées de }\delta\Phi,\text{ qui ne}\\
&\qquad\text{sont pas }\in\Phi.
\end{aligned}
\]

% source: middle, redundant vertices
On aura
\[
\delta\Phi^*=\delta\Phi\setminus\Phi_{\deleted{\unsure{\mathrm{le}}}\mathrm{red}},\quad\text{soit }T'
\]
où $\Phi_{\mathrm{réd}}$ est l'ens. des sommets redondants de $\Phi$, i.e.
\[
\Phi_{\mathrm{red}}=\{s\in\Phi_{\unsure{0|o}}\mid\operatorname{ord}(s\unsure{.|,}\Phi)=2\}
\]

% source: middle-lower, digression
\deleted{On appelle}

Il y a lieu ici de faire la petite digression sur la nature des sommets d'un $\unsure{\widehat M|\widehat{\mathscr M}}$-polygone (cinq types: sommets isolés, sommets-bord propres et impropres, sommets internes propres (ou redondants) ou impropres (dits lacunaires)) -- cf. pages \unsure{110|10}, \unsure{111|11}. On trouve \deleted{\unsure{aussitôt}} que $\Phi^*$ n'a pas de sommets redondants, donc $\delta\Phi^{\unsure{\circ *|**}}=\delta\Phi^*$.

% source: bottom, double complement
On obtient
\[
\begin{aligned}
\deleted{\unsure{\delta}}\Phi^{**}&=\text{ens des strates interstitielles}\inserted{X}\text{ de }T'=\delta\Phi\setminus\Phi_{\mathrm{red}}\\
&\quad\deleted{\text{\unsure{pour lesquels il existe}}}\inserted{\text{telles qu'il existe}}Y\in\Phi\ \unsure{\text{avec}|\text{tel que}}\ Y\ll X.
\end{aligned}
\]
On \unsure{sait|voit} que $\Phi^{\unsure{**}}$ s'obtient à partir de $\Phi$
% archive sheet number: 13
\scanpage{15}
\mspage{11}
% source: top, continuation from scan 14
par ``effacement des sommets redondants'', \deleted{et}\unsure{ou} encore par ``composition des \deleted{\unsure{segments}} intervalles adjacents aux sommets redondants''. On a
\[
({\unsure{\Phi|\varphi}}^{\unsure{**|\circ *}})_{\unsure{0|o}}=\Phi_0\setminus\Phi_{\mathrm{red}}
\]
\[
\begin{aligned}
(\Phi_{\unsure{c|0|\epsilon|\cdot}}^{**})_1&=\deleted{\text{composé}}\text{ ens des intervalles}\\
&\qquad\deleted{\text{composés}}\inserted{\text{obtenus}}\text{ par composition des}\\
&\qquad\deleted{\text{\unsure{suite}}}\text{drapeaux maximaux}\\
&\qquad\uline{\text{connexes}}\text{ contenus dans }\Phi,
\end{aligned}
\]
% small c-like mark at lower right of Phi in the second display retained above; its status is uncertain
% source: middle, chain and relations
i.e. de la forme
\[
X_{\unsure{1|l}},x_1,X_2,x_2,\cdots,x_{\unsure{k|n}-1},X_k
\]
avec
\[
\deleted{x_1\triangleleft X_1,X_2,\ \unsure{x}}
\]
\[
\begin{aligned}
&X_1<x_1<X_2<x_2\ \cdots\ <x_{\unsure{n|k|u}-1}<X_n\\
&\text{et }x_1\triangleleft X_1,X_2,\quad x_2\unsure{\mid|,}\triangleleft X_2,X_3,\quad\cdots,\quad x_{n-1}\triangleleft X_{n-1},X_n.
\end{aligned}
\]

% source: middle-lower, components
\unsure{Elles|Ils} sont en corr. 1-1 avec les composantes connexes de $\Phi$, lesquelles n'en diffèrent que par l'adjonction éventuelle d'un $x_0\triangleleft X_1$, ou d'un $x_n\triangleleft X_n$, ou des deux.

% source: lower, shadow of a polygon
\deleted{\unsure{Et}} Ceci posé, à chaque \uline{polygone} $\Phi$, associons
\[
\operatorname{Omb}^{\unsure{\circ|0}}_{\unsure{M|\mathscr M}}(\Phi)=\deleted{\unsure{\forall X\in}}\bigcup_{X\in\Phi}\operatorname{Omb}^{\circ}_M(X)
=\left\{Z\in M\left|\begin{aligned}&\exists X\in\Phi,\\&Z\overset{\circ}{\ll}X\end{aligned}\right.\right\}
\]

% source: bottom, introduction to following proposition
Ceci posé, le résultat \unsure{\uline{clef}|clef} est le suivant
% archive sheet number: 14
\scanpage{16}
\mspage{12}
% source: top, theorem and revised hypotheses
\uline{Th} Soit $\unsure{\Phi|\varphi}$ un $\unsure{\widehat M|\widehat{\mathscr M}}$-polygone,\inserted{non vide,} \deleted{et supposons que \unsure{a)} $\operatorname{or}\delta\Phi$ \unsure{en soit a} \unsure{a)} $\operatorname{or}\delta\Phi$.}

% source: upper part, entire hypothesis block cancelled by oblique strokes
\deleted{soit $a=\operatorname{or}\delta\Phi$, $b=\operatorname{ex}\delta\Phi$, et supposons que}

\deleted{a) Si \inserted{Si $\unsure{L_{<a}|I_{<a}}\in\Phi$ et,} $a$ est minimal, i.e un plus petit élément, et}

\deleted{b) Si \inserted{Si $L_{>b}\in\Phi$ et,} $b$ est maximal, i.e un plus grand élément}

\deleted{(Ce qui résulte\inserted{\unsure{aussi}} de l'hypothèse \unsure{Hyp}: $\unsure{\forall|L}$ \deleted{\unsure{Si}} $\unsure{L|\mathscr L}$ a un plus petit élément, ou bien pour $\forall a\in L$, $L_{<a}\ne\varnothing$, et $L$ a un plus grand élément, ou pour $\forall a\in L$, $L_{>a}\ne\varnothing$.)}

% source: middle, theorem conclusion
Alors
\[
\mathrm{cos}\!\!\unsure{-}\!{\unsure{\psi|w|W}}^{\unsure{\circ|0}}_M(\Phi)=\mathrm{cos}\!\!-\!\psi^{\circ}(\operatorname{Omb}^{\circ}_M(\Phi))=\operatorname{Omb}^{\circ}_M(\Phi^*)
\]

% source: lower-middle, first corollary
\uline{Corollaire 1} Sous la condition Hyp, pour simplifier:
\[
\mathrm{s}\!-\!\psi^{\circ}_M(\Phi)=\mathrm{s}\!-\!\psi^{\circ}_M(\operatorname{Omb}^{\circ}_M(\Phi))=\operatorname{Omb}^{\circ}_M(\Phi^{**})
\]

% source: lower text, second corollary begins
\uline{Corollaire 2} \deleted{Soit (\unsure{sous la condition}} Soit $T\in\operatorname{Drap}(L)$. Considérons l'ens des figures telles que
% archive sheet number: 15
\scanpage{17}
\mspage{13}
% source: top, date and headings
6. juillet.

\unsure{Re-venons|Revenons}, sur l'introduction des $I_{<a},I_{>b}$.

Autre présentation\unsure{:|.}

% source: upper part, ordered segment
$1^\circ$) Pour y comprendre quelque chose, \unsure{commencons|commençons} avec le cas d'un ``\deleted{ensemble}\inserted{\uline{segment}} \uline{ordonné}'' ou ``\uline{segment orienté combinatoire}'' (peut-être faudrait-il les appeler ``\uline{pseudo-segment}'', car on sera amené par la suite à restreindre les \unsure{hypothèses|hypothéses}), c'est le \uline{système}
\[
\{I,\omega_{\unsure{0|o}},\omega_{\unsure{1|l}}\}\quad\text{ou}\quad(I,\unsure{\partial|\delta}I),\quad\partial I=\{\omega_0,\omega_1\}
\]
d'un ensemble ordonné $I$\inserted{(ou \uline{support du segment orienté})}, d'un plus petit élément $\omega_0$ ou \uline{origine}, et d'un plus grand élément $\omega_1$ ou \uline{extrémité} avec
\[
\omega_0<\omega_1\qquad\text{i.e.}\quad\operatorname{card}I\geq2.
\]

% source: middle, two-element case
On n'exclut pas le cas de $\operatorname{card}I=2$, i.e. $\unsure{\lvert I\rvert^{\circ}|\mathbb I^{\circ}|\lvert I\rvert^0|\mathbb I^0}=\varnothing$, où
\[
\lvert I\rvert^{\circ}=I\setminus\partial I\subset I.
\]
% source: angled left margin
% slanted
\margin{\unsure{C'est le seul} cas où on trouve une maquette}
Ce cas est celui des \uline{segments indivisibles}, qui donnent lieu à une \unsure{m.g. ori} \unsure{qui est} \unsure{une} \uuline{maquette} simpliciale, \inserted{\unsure{savoir},} celle du simplexe orienté-type de dim 1, noté $\Delta^1$.

% source: lower part, restriction to finite intervals
Dans ce cas, on n'introduit pas des strates de la forme $I_{<a},I_{>b}$. La notion du polygone \unsure{complémentaire} ${\unsure{\Phi|\varphi}}^*$ d'un polygone $\Phi$ se développe, en restant dans le cadre des $a$ ou $P_a$ ($a\in I$) et $I_{a,b}$ ($a<b$, $a,b\in I$).

Les développements concernant la détermination
% archive sheet number: 16
